% ============================================================================
% LANGENTWURF LE-04b: El horario (Stundenplan, Wochentage, Uhrzeiten)
% Spanisch Klasse 7 - Sequenz 4: Mi colegio
% ============================================================================
% Tier: 1 (vollständig)
% Doppelstunde: 4b (Std. 3-4 von 8)
% Fachfarbe: #c41e3a (Karminrot)
% ============================================================================

\documentclass[11pt,a4paper]{article}

\usepackage[utf8]{inputenc}
\usepackage[T1]{fontenc}
\usepackage[ngerman]{babel}
\usepackage{geometry}
\usepackage{fancyhdr}
\usepackage{lastpage}
\usepackage{xcolor}
\usepackage{tcolorbox}
\usepackage{enumitem}
\usepackage{tabularx}
\usepackage{longtable}
\usepackage{array}
\usepackage{booktabs}
\usepackage{amssymb}

% ============================================================================
% KONFIGURATION
% ============================================================================

\newcommand{\fach}{Spanisch}
\newcommand{\fachfarbe}{c41e3a}
\newcommand{\jahrgangsstufe}{7}
\newcommand{\schulart}{Gesamtschule}
\newcommand{\stundenthema}{El horario -- Stundenplan und Uhrzeiten}
\newcommand{\reihenthema}{Mi colegio}
\newcommand{\stundennummer}{4b}
\newcommand{\reihenstunden}{4}
\newcommand{\gession}{A1}
\newcommand{\lehrwerk}{Qué pasa 1}
\newcommand{\lehrwerkseite}{S. 48-53}
\newcommand{\lerngruppengroesse}{24}

% ============================================================================
% LAYOUT
% ============================================================================
\geometry{left=2.5cm, right=2.5cm, top=2.5cm, bottom=2.5cm}
\definecolor{fach}{HTML}{\fachfarbe}
\definecolor{gray}{HTML}{666666}
\definecolor{lightgray}{HTML}{f5f5f5}
\definecolor{successgreen}{HTML}{2a7a4b}
\definecolor{warningorange}{HTML}{b35c00}
\definecolor{basis}{HTML}{2a7a4b}
\definecolor{standard}{HTML}{2c5aa0}
\definecolor{erweiterung}{HTML}{7b1fa2}

\tcbuselibrary{skins,breakable}

\newtcolorbox{infobox}[1][]{
    colback=lightgray,
    colframe=fach,
    fonttitle=\bfseries,
    title=#1,
    breakable
}

\newtcolorbox{scorebox}{
    colback=white,
    colframe=fach,
    fonttitle=\bfseries,
    title=Didaktik-Score,
    breakable
}

\pagestyle{fancy}
\fancyhf{}
\fancyhead[L]{\small\textcolor{gray}{\fach{} -- Jg. \jahrgangsstufe{} -- \stundenthema}}
\fancyhead[R]{\small\textcolor{gray}{LE-\stundennummer{} | Tier 1}}
\fancyfoot[C]{\small\thepage{} / \pageref{LastPage}}
\renewcommand{\headrulewidth}{0.4pt}

% Differenzierungsmarker
\newcommand{\B}{\textcolor{basis}{\textbf{[B]}}}
\newcommand{\Std}{\textcolor{standard}{\textbf{[S]}}}
\newcommand{\E}{\textcolor{erweiterung}{\textbf{[E]}}}

% ============================================================================
% DOKUMENT
% ============================================================================
\begin{document}

% ============================================================================
% DECKBLATT
% ============================================================================
\begin{titlepage}
    \centering
    \vspace*{2cm}
    {\large\textcolor{gray}{\schulart}}\\[2cm]
    {\Huge\textcolor{fach}{\textbf{Langentwurf}}}\\[0.5cm]
    {\large LE-\stundennummer{} | Tier 1 (vollständig)}\\[2cm]
    {\LARGE\textbf{\stundenthema}}\\[0.5cm]
    {\large Sequenz: \reihenthema}\\[2cm]
    \begin{tabular}{rl}
        \textbf{Fach:} & \fach{} (A1) \\[0.3cm]
        \textbf{Klasse:} & \jahrgangsstufe \\[0.3cm]
        \textbf{Lehrwerk:} & \lehrwerk, \lehrwerkseite \\[0.3cm]
        \textbf{Doppelstunde:} & \stundennummer{} (Std. 3-4 von 8) \\[0.3cm]
    \end{tabular}
    \vfill
\end{titlepage}

\tableofcontents
\newpage

% ============================================================================
% 1. BEDINGUNGSANALYSE
% ============================================================================
\section{Bedingungsanalyse}

\subsection{Lerngruppe}
\begin{tabular}{ll}
    Kursgröße: & \lerngruppengroesse{} SuS \\
    Jahrgangsstufe: & \jahrgangsstufe \\
    Sprachniveau: & A1 (GER) -- Anfänger \\
    Fremdsprache: & 2. Fremdsprache (nach Englisch) \\
\end{tabular}

\subsection{Vorwissen aus 04a}
\begin{itemize}
    \item Schulfächer (\textit{las asignaturas}): matemáticas, inglés, español, educación física, etc.
    \item Verb \textit{tener}: yo tengo, tú tienes, él/ella tiene, nosotros tenemos
    \item Frage: ¿Qué asignaturas tienes?
    \item Redemittel: Tengo matemáticas, Me gusta/No me gusta...
\end{itemize}

\subsection{Differenzierung (intern)}
\begin{tabular}{|l|p{3.5cm}|p{3.5cm}|p{3.5cm}|}
\hline
& \B{} \textbf{Basis} & \Std{} \textbf{Standard} & \E{} \textbf{Erweiterung} \\
\hline
\textbf{Ziel} & BR: Rezeption & MR: Produktion mit Hilfe & GY: Freie Produktion \\
\hline
\textbf{Wochentage} & Nachsprechen, zuordnen & Selbstständig benennen & In Sätzen verwenden \\
\hline
\textbf{Uhrzeiten} & Volle Stunden & + halbe Stunden & + Viertelstunden \\
\hline
\textbf{Dialog} & Mit Vorlage & Mit Stichworten & Freier Dialog \\
\hline
\end{tabular}

% ============================================================================
% 2. SACHANALYSE
% ============================================================================
\section{Sachanalyse}

\subsection{Sprachliche Strukturen}

\textbf{Die Wochentage (\textit{los días de la semana}):}
\begin{center}
\begin{tabular}{|l|l|l|}
\hline
\textbf{Spanisch} & \textbf{Deutsch} & \textbf{Abkürzung} \\
\hline
\textbf{el lunes} & Montag & L \\
\textbf{el martes} & Dienstag & M \\
\textbf{el miércoles} & Mittwoch & X \\
\textbf{el jueves} & Donnerstag & J \\
\textbf{el viernes} & Freitag & V \\
\textbf{el sábado} & Samstag & S \\
\textbf{el domingo} & Sonntag & D \\
\hline
\end{tabular}
\end{center}

\textbf{Hinweise:}
\begin{itemize}
    \item Wochentage sind maskulin und werden kleingeschrieben
    \item Miércoles hat einen Akzent (é)
    \item Mittwoch wird mit X abgekürzt (M = martes)
\end{itemize}

\textbf{Uhrzeiten (\textit{las horas}):}
\begin{center}
\begin{tabular}{|l|l|}
\hline
\textbf{Spanisch} & \textbf{Bedeutung} \\
\hline
\textit{¿Qué hora es?} & Wie spät ist es? \\
\textit{Es la una.} & Es ist ein Uhr. \\
\textit{Son las dos / tres / ...} & Es ist zwei / drei / ... Uhr. \\
\textit{Son las ocho y media.} & Es ist halb neun. \\
\textit{Son las nueve y cuarto.} & Es ist Viertel nach neun. \\
\textit{Son las diez menos cuarto.} & Es ist Viertel vor zehn. \\
\hline
\end{tabular}
\end{center}

\textbf{Zeitangaben mit Präpositionen:}
\begin{itemize}
    \item \textit{a las ocho} = um acht Uhr
    \item \textit{de ocho a nueve} = von acht bis neun
    \item \textit{los lunes} = montags (regelmäßig)
\end{itemize}

\textbf{Fragen zum Stundenplan:}
\begin{itemize}
    \item \textit{¿Qué día es hoy?} = Welcher Tag ist heute?
    \item \textit{¿A qué hora tienes matemáticas?} = Um wie viel Uhr hast du Mathe?
    \item \textit{¿Cuándo tienes español?} = Wann hast du Spanisch?
\end{itemize}

\subsection{Kontrastive Betrachtung}
\begin{itemize}
    \item \textbf{Positiv:} Struktur ähnlich Englisch (What time is it? / What day is it?)
    \item \textbf{Unterschied:} ``halb neun'' = \textit{ocho y media} (8 + 30, nicht 9 - 30!)
    \item \textbf{Interferenz:} Deutsche Schüler sagen oft *\textit{media nueve} statt \textit{ocho y media}
    \item \textbf{Grammatik:} \textit{Es la una} (Singular) vs. \textit{Son las dos} (Plural)
\end{itemize}

% ============================================================================
% 3. DIDAKTISCHE ANALYSE
% ============================================================================
\section{Didaktische Analyse}

\subsection{Stellung in der Sequenz}
Dies ist die \textbf{2. Doppelstunde} (Std. 3-4 von 8) der Sequenz ``\reihenthema''.

\subsection{Kompetenzziele (Rahmenplan MV)}
\begin{itemize}
    \item \textbf{Kommunikativ:} Nach Uhrzeiten und Wochentagen fragen und antworten
    \item \textbf{Sprachlich:} Wochentage und Uhrzeiten korrekt verwenden
    \item \textbf{Methodisch:} Einen Stundenplan auf Spanisch lesen und beschreiben
\end{itemize}

\subsection{Legitimation}
\textbf{Rahmenplan Spanisch MV:}
\begin{itemize}
    \item Sprechen: Alltagssituationen bewältigen (Uhrzeit, Termine)
    \item Verfügen über sprachliche Mittel: Wortschatz -- Schule, Tagesablauf
\end{itemize}

% ============================================================================
% 4. STUNDENZIELE
% ============================================================================
\section{Stundenziele}

\subsection{Grobziel}
Die SuS erweitern ihre \textbf{kommunikative Kompetenz}, indem sie Wochentage und Uhrzeiten benennen und ihren Stundenplan auf Spanisch beschreiben können.

\subsection{Feinziele (operationalisiert)}
\begin{tabular}{|c|p{8cm}|c|c|}
\hline
\textbf{Nr.} & \textbf{Die SuS können...} & \textbf{AFB} & \textbf{Diff.} \\
\hline
FZ1 & die sieben Wochentage auf Spanisch benennen und in die richtige Reihenfolge bringen & I & alle \\
\hline
FZ2 & Uhrzeiten (volle und halbe Stunden) verstehen und angeben & I/II & alle \\
\hline
FZ3 & nach dem Tag und der Uhrzeit fragen: ¿Qué día es hoy? ¿A qué hora...? & II & alle \\
\hline
FZ4 & einen Dialog über den Stundenplan führen: ¿Cuándo tienes...? & II/III & \Std{}/\E{} \\
\hline
\end{tabular}

% ============================================================================
% 5. METHODISCHE ANALYSE
% ============================================================================
\section{Methodische Analyse}

\subsection{Medien}
\begin{itemize}
    \item Tafel/Whiteboard
    \item Lehrwerk: \lehrwerk, \lehrwerkseite
    \item Arbeitsblatt: AB-04b.pdf
    \item Bildkarten: Wochentage, Uhren
    \item Audio: Track (Wochentage-Lied)
    \item Optional: LP-04b.html (digital)
\end{itemize}

\subsection{Sozialformen}
\begin{itemize}
    \item Plenum: Einführung Wochentage, Uhrzeiten
    \item Partnerarbeit: Stundenplan-Dialog
    \item Einzelarbeit: AB-Aufgaben
\end{itemize}

% ============================================================================
% 6. VERLAUFSPLANUNG
% ============================================================================
\section{Verlaufsplanung}

\textbf{1. Stunde (45 Min.) -- Schwerpunkt: Wochentage + Uhrzeiten}

\begin{longtable}{|p{1cm}|p{1.5cm}|p{4cm}|p{3cm}|p{2cm}|p{2cm}|}
\hline
\textbf{Zeit} & \textbf{Phase} & \textbf{Lehrerhandeln} & \textbf{Schülerhandeln} & \textbf{SF} & \textbf{Medien} \\
\hline
\endhead

0--5 & Einstieg & Kalender zeigen: ``¿Qué día es hoy?'' & Vermutungen äußern & Plenum & Kalender \\
\hline

5--15 & Input 1 & Wochentage einführen, vorsprechen, Tafelanschrieb & Nachsprechen, Notieren & Plenum & Tafel \\
\hline

15--20 & Übung 1 & Wochentage-Memory / Zuordnung & Spielen, üben & PA & Karten \\
\hline

20--30 & Input 2 & Uhrzeiten einführen mit Uhr: ``Son las...'' & Nachsprechen, Mitmachen & Plenum & Uhr \\
\hline

30--40 & Übung 2 & AB-04b Aufg. 1+2: Wochentage + Uhrzeiten & Schriftlich bearbeiten & EA & AB \\
\hline

40--45 & Sicherung & Lösungen besprechen, Fragen klären & Vergleichen & Plenum & Tafel \\
\hline
\end{longtable}

\vspace{1em}

\textbf{2. Stunde (45 Min.) -- Schwerpunkt: Stundenplan + Dialog}

\begin{longtable}{|p{1cm}|p{1.5cm}|p{4cm}|p{3cm}|p{2cm}|p{2cm}|}
\hline
\textbf{Zeit} & \textbf{Phase} & \textbf{Lehrerhandeln} & \textbf{Schülerhandeln} & \textbf{SF} & \textbf{Medien} \\
\hline
\endhead

0--5 & Aktivierung & Schnelles Uhrzeiten-Quiz & Antworten & Plenum & Uhr \\
\hline

5--15 & Input & Stundenplan vorstellen, Fragen einführen: ¿A qué hora...? ¿Cuándo...? & Mitschreiben, Nachsprechen & Plenum & Tafel \\
\hline

15--25 & Übung 1 & AB-04b Aufg. 3: Stundenplan ausfüllen & Eintragen & EA & AB \\
\hline

25--35 & Übung 2 & PA: Dialog zum Stundenplan ``¿Cuándo tienes...?'' & Partner befragen & PA & -- \\
\hline

35--40 & Sicherung & AB-04b Aufg. 4: Dialog aufschreiben & Schreiben & EA & AB \\
\hline

40--45 & Abschluss & Zusammenfassung, HA erklären & Aufschreiben & Plenum & -- \\
\hline
\end{longtable}

\textbf{Legende:} EA = Einzelarbeit, PA = Partnerarbeit, SF = Sozialform

% ============================================================================
% 7. ERWARTETE SCHWIERIGKEITEN
% ============================================================================
\section{Erwartete Schwierigkeiten}

\begin{tabular}{|p{5cm}|p{8cm}|}
\hline
\textbf{Schwierigkeit} & \textbf{Reaktion} \\
\hline
``Halb neun'' vs. \textit{ocho y media} & Visualisierung: Uhr zeigen, ``y media = +30'' \\
\hline
Miércoles-Aussprache & Drill: mi-ér-co-les, Akzent betonen \\
\hline
\textit{Es la una} vs. \textit{Son las...} & Regel: 1 Uhr = Singular, Rest = Plural \\
\hline
Wochentage durcheinander & Merkspruch/Lied, häufige Wiederholung \\
\hline
Präposition \textit{a las} vergessen & Formel: Uhrzeit $\rightarrow$ immer \textit{a las} + Zahl \\
\hline
\end{tabular}

% ============================================================================
% 8. HAUSAUFGABE
% ============================================================================
\section{Hausaufgabe}

\begin{itemize}
    \item Lehrwerk S. 50, Übung 6 (Uhrzeiten schreiben)
    \item Vokabeln lernen: 7 Wochentage + Uhrzeiten
    \item Optional: LP-04b.html bearbeiten (interaktive Übungen)
\end{itemize}

% ============================================================================
% 9. LÖSUNGEN
% ============================================================================
\section{Lösungen AB-04b}

\textbf{Merkkasten Wochentage:}
\begin{itemize}
    \item lunes, martes, miércoles, jueves, viernes, sábado, domingo
\end{itemize}

\textbf{Merkkasten Uhrzeiten:}
\begin{itemize}
    \item \textit{Son las ocho.} = Es ist acht Uhr.
    \item \textit{Son las nueve y media.} = Es ist halb zehn.
    \item \textit{a las diez} = um zehn Uhr
\end{itemize}

\textbf{Aufgabe 1 (Wochentage ordnen):} (4P)
\begin{enumerate}
    \item lunes
    \item martes
    \item miércoles
    \item jueves
    \item viernes
    \item sábado
    \item domingo
\end{enumerate}

\textbf{Aufgabe 2 (Uhrzeiten schreiben):} (4P)
\begin{enumerate}[label=\alph*)]
    \item Son las ocho.
    \item Son las diez y media.
    \item Es la una.
    \item Son las tres y media.
\end{enumerate}

\textbf{Aufgabe 3 (Stundenplan ausfüllen):} (6P)
\begin{itemize}
    \item Individuelle Lösung je nach Stundenplan
    \item Bewertung: Korrekte Fächer (3P) + korrekte Zeiten (3P)
\end{itemize}

\textbf{Aufgabe 4 (Dialog):} (6P)
\begin{itemize}
    \item A: ¿Cuándo tienes matemáticas?
    \item B: Los lunes a las ocho.
    \item A: ¿Y español?
    \item B: Los martes a las diez.
    \item A: ¿A qué hora tienes educación física?
    \item B: Los viernes a las once.
\end{itemize}

% ============================================================================
% 10. DIDAKTIK-SCORE
% ============================================================================
\newpage
\section{Didaktik-Score}

\begin{scorebox}
\textbf{Bewertung der didaktischen Qualität nach 10 Dimensionen}\\
\textbf{Ziel-Score: $\geq$ 9.0 (Premium-Qualität)}

\vspace{1em}
\begin{tabular}{|c|l|c|c|p{4.5cm}|}
\hline
\textbf{\#} & \textbf{Dimension} & \textbf{Gew.} & \textbf{Score} & \textbf{Notiz} \\
\hline
1 & Differenzierung & 15\% & 9/10 & [B] volle Std., [S] +halbe, [E] +Viertel \\
\hline
2 & Taxonomie (AFB) & 10\% & 9/10 & AFB I: 35\%, AFB II: 50\%, AFB III: 15\% \\
\hline
3 & Lernziel-Alignment & 10\% & 10/10 & Rahmenplan: Uhrzeit, Alltagssprache \\
\hline
4 & Aktivierung & 10\% & 9/10 & PA-Dialog, Memory, Quiz \\
\hline
5 & Scaffolding & 10\% & 9/10 & Volle Std. $\rightarrow$ halbe $\rightarrow$ Dialog \\
\hline
6 & Feedback-Qualität & 10\% & 9/10 & Elaboriertes Feedback in LP \\
\hline
7 & Sprachliche Klarheit & 10\% & 10/10 & Altersgerecht, kurze Sätze \\
\hline
8 & Motivation \& Relevanz & 10\% & 10/10 & Eigener Stundenplan = Alltagsrelevanz \\
\hline
9 & Inklusion (UDL) & 10\% & 9/10 & Visuell (Uhr) + auditiv + kinästhetisch \\
\hline
10 & Praktikabilität & 5\% & 9/10 & 2x45 Min realistisch \\
\hline
\multicolumn{3}{|r|}{\textbf{GESAMT:}} & \textbf{9.3/10} & \textcolor{successgreen}{\textbf{FREIGABE}} \\
\hline
\end{tabular}
\end{scorebox}

\end{document}
